\documentclass[10pt,a4paper]{amsart}
\usepackage[margin=19mm]{geometry}
\usepackage{amsmath,amssymb,mathtools}
\usepackage[scr=rsfs,cal=euler]{mathalfa}
\usepackage{xfrac}
\usepackage[unicode,hidelinks,bookmarks=false]{hyperref}
\newcommand{\ie}{i.e.\ }
\newcommand{\cf}{cf.\ }
\newcommand{\resp}{respectively\ }
\newcommand{\Subsection}{\subsection}
\providecommand{\nobreakdash}{\nobreak\hbox{-}}
\setlength{\emergencystretch}{2em}
\multlinegap0pt
\usepackage{bm}
\usepackage{bbm}
\usepackage{dsfont}
\usepackage{xspace}
\usepackage{cancel}
\usepackage{mathtools}
\usepackage{amsthm}
\makeatletter
\@ifundefined{theorem}{\newtheorem{theorem}{Theorem}}{}
\@ifundefined{assumption}{\newtheorem{assumption}[theorem]{Assumption}}{}
\@ifundefined{proposition}{\newtheorem{proposition}{Proposition}}{}
\makeatother
\newcommand{\cactive}{\mathcal{B}}
\newcommand{\cinactive}{\mathcal{N}}
\newcommand{\epstol}{\mathbf{u}}
\title[Condensed Interior-Point Methods for Scalable Nonlinear Prog]{Condensed Interior-Point Methods for Scalable Nonlinear Programming on GPUs}
\author{Fran\c{c}ois Pacaud, Sungho Shin, Alexis Montoison, Michel Schanen, Mihai Anitescu}
\date{Source: arXiv:2405.14236v3 (CC BY 4.0)}
\begin{document}
\maketitle

\noindent\emph{Source and licence.} Condensed Interior-Point Methods for Scalable Nonlinear Programming on GPUs. Version arXiv:2405.14236v3 (CC BY 4.0), distributed under the Creative Commons Attribution 4.0 International licence, as recorded in the arXiv licence metadata for this exact version and shown on the abstract page https://arxiv.org/abs/2405.14236v3. Full rights notice: licenses/arxiv-2405.14236-CC-BY-4.0.txt.\par\smallskip

\noindent\textit{Editorial scope.} Counted unit: the Proposition whose source label is \texttt{eq:cond:errorrgamma} in the article (source0.tex lines 1202--1208, source label \texttt{eq:cond:errorrgamma}), delivered together with the article's own complete original proof of it (source0.tex lines 1209--1257, one \texttt{proof} environment of 49 lines). No mathematics is written, repaired or re-derived: statement and proof are the article's own text line for line. Required context retained uncounted: eq:centralitycond (lines 744--751); eq:cond:invariantsubpsace (lines 923--927); eq:cond:boundinvariantsubspace (lines 943--951); eq:cond:backwardstable (lines 1031--1034); hyp:cond:wellcond (lines 1041--1051); eq:cond:inversecondensed (lines 1084--1094); eq:cond:rhserrorfull (lines 1120--1125); eq:cond:invperturbed (lines 1129--1135); eq:cond:condensedrhs (lines 1165--1173). Every definition, display and label of those lines is retained unchanged. Omitted from this package and not redistributed: 1--743, 752--922, 928--942, 952--1030, 1035--1040, 1052--1083, 1095--1119, 1126--1128, 1136--1164, 1174--1201, 1258--2231. Nothing inside a retained excerpt is omitted or altered: every retained body is delivered exactly as the article's lines are written, so no omission receipt of any kind accompanies this package. Citations: the retained excerpts carry no citation command at all, so no bibliography entry of the article is transcribed and no reference is redistributed. Reference adaptation: none was needed and none was made; every reference inside the delivered unit resolves inside it, so no unresolved reference or citation is produced. Printed slips kept literal: no printed word, symbol, hypothesis or number of the counted statement or of its proof was altered. Layout and class: the article's own class is \texttt{svjour3} with packages including amsmath, amsfonts, amssymb, xcolor, bm, url, booktabs, array, tikz, xspace, hyperref; the wrapper is the lane's pinned \texttt{amsart} editorial wrapper at 10pt on A4, which restates the theorem-like environments the retained text needs and adds the article's own macro definitions that the retained text needs (\texttt{\textbackslash cactive}, \texttt{\textbackslash cinactive}, \texttt{\textbackslash epstol}), with extra wrapper packages (bm, bbm, dsfont, xspace, cancel, mathtools). The delivered numbering is the unit's own. Assets: no retained span contains a figure environment, a graphics-inclusion command, a PDF-page inclusion, a table or a TikZ picture; no image file is copied into this package, the package declares no asset dependency and no image file is redistributed with it. Licence: version arXiv:2405.14236v3 of this article is distributed under the Creative Commons Attribution 4.0 International licence, as recorded in the arXiv licence metadata for this exact version and shown on the abstract page https://arxiv.org/abs/2405.14236v3. A full rights notice is delivered at licenses/arxiv-2405.14236-CC-BY-4.0.txt. Characters: every retained excerpt is pure ASCII, so the delivered engine, XeLaTeX with its default Latin Modern text font, reports no missing character.
\begin{subequations}
  \label{eq:centralitycond}
  \begin{align}
    & \| \nabla_p \mathcal{L}(p, v) \| \leq C \; \Xi(s, v) \;,  \\
    \label{eq:centralitycond:complement}
    & (s, v) > 0 \;,\quad s_i v_i \geq \alpha \, \Xi(s, v) \quad \forall i =1, \cdots, m_i \; ,
  \end{align}
\end{subequations}

  \begin{equation}
    \label{eq:cond:invariantsubpsace}
    U_L - Y = O(\underline{\sigma}^{-1}) \; , \quad
    U_S - Z = O(\underline{\sigma}^{-1}) \; .
  \end{equation}

\begin{equation}
  \label{eq:cond:boundinvariantsubspace}
  \begin{aligned}
    & \| K_\gamma \| = \| \Sigma_L \| = \Theta(\overline{\sigma}) \; , &
    \Sigma_L^{-1} = O(\underline{\sigma}^{-1})  \;, \\
    & \| K_\gamma^{-1} \| = \| \Sigma_S^{-1} \| = \Theta(1) \, , &
  \Sigma_S = \Theta(1) \, .
  \end{aligned}
\end{equation}

\begin{equation}
  \label{eq:cond:backwardstable}
  \|\Delta_s K_\gamma\| \leq \epstol \varepsilon_n \|K_\gamma\| \;,
\end{equation}

\begin{assumption} Let $(p, v)$ be the current primal-dual iterate. We assume:
  \begin{itemize}
    \item[(a)] $(p, v)$ satisfies the centrality conditions~\eqref{eq:centralitycond}.
    \item[(b)] The parameter $\gamma$ satisfies $\gamma = \Theta(\Xi^{-1})$.
    \item[(c)] The duality measure is large enough relative to the precision $\epstol$: $\epstol \ll \Xi$.
    \item[(d)] The primal step $\widehat{x}$ is computed using a backward
      stable method satisfying~\eqref{eq:cond:backwardstable} for a small constant
      $\varepsilon_n$.
  \end{itemize}
  \label{hyp:cond:wellcond}
\end{assumption}

\begin{equation}
  \label{eq:cond:inversecondensed}
  K_\gamma^{-1}  = \begin{bmatrix} U_L & U_S \end{bmatrix}
  \begin{bmatrix}
    \Sigma_L^{-1} & 0 \\ 0 & \Sigma_S^{-1}
  \end{bmatrix}
  \begin{bmatrix}
    U_L^\top \\ U_S^\top
  \end{bmatrix}
  \; .
\end{equation}

\begin{equation}
  \label{eq:cond:rhserrorfull}
  \| \Delta x \|
  \leq  \| \Sigma_L^{-1} \| \| \Delta b_L \| + \| \Sigma_S^{-1} \| \| \Delta b_S \| =
  O(\|\Delta b \|) \; .
\end{equation}

\begin{equation}
  \label{eq:cond:invperturbed}
  \begin{aligned}
    (K_\gamma + \Delta K_\gamma)^{-1} &= (I + K_\gamma^{-1} \Delta K_\gamma)^{-1} K_\gamma^{-1} \; , \\
                                      &= K_\gamma^{-1} - K_\gamma^{-1}\Delta K_\gamma K_\gamma^{-1} + O(\|\Delta K_\gamma\|^2) \; .
  \end{aligned}
\end{equation}

\begin{equation}
  \label{eq:cond:condensedrhs}
  \left\{
  \begin{aligned}
    \bar{r}_1 &= - \widehat{r}_1 + \widehat{H}^\top\big(\widehat{D}_{s} \widehat{r}_{4} - \widehat{r}_{2} \big) \;, \\
     \bar{r}_2 &= -\widehat{r}_3 \; .
  \end{aligned}
  \right.
\end{equation}

\begin{proposition}
In floating point arithmetic, the error in the right-hand-side $\Delta \widehat{r}_\gamma$ satisfies:
\begin{equation}
  \label{eq:cond:errorrgamma}
  \Delta \widehat{r}_\gamma = -\Delta \bar{r}_2 + \widehat{G} \widehat{K}_\gamma^{-1} \Delta s_\gamma = O(\epstol) \;.
\end{equation}
\end{proposition}

\begin{proof}
Using \eqref{eq:cond:condensedrhs}, we have
\begin{equation*}
  \label{eq:cond:boundderivationhykkt}
  \begin{aligned}
  \bar{r}_1 + \gamma \widehat{G}^\top \bar{r}_2 &=
- \widehat{r}_1 + \gamma \widehat{G}^\top \widehat{r}_3+ \widehat{H}^\top\big(\widehat{D}_{s} \widehat{r}_{4} - \widehat{r}_{2} \big) \\
&=  -
\underbrace{\widehat{r}_1}_{O(\epstol)} +
\underbrace{
\widehat{H}_{\cinactive}^\top\big(\widehat{D}_{s,\cinactive} \widehat{r}_{4,\cinactive} - \widehat{r}_{2,\cinactive} \big)}_{O(\Xi \epstol)}
+ \underbrace{\widehat{A}^\top \begin{bmatrix}
  \widehat{D}_{s,\cactive} \widehat{r}_{4,\cactive} - \widehat{r}_{2,\cactive}  \\
  \gamma \widehat{r}_3
\end{bmatrix}}_{O(\Xi^{-1}\epstol)} \; .
  \end{aligned}
\end{equation*}
The error decomposes as $\Delta s_\gamma = Y \Delta s_Y  + Z \Delta s_Z
= U_L \Delta s_L + U_S \Delta s_S$.
We have $\Delta s_Y = O(\Xi^{-1} \epstol)$ and $\Delta s_Z = O(\epstol)$.
Using \eqref{eq:cond:invariantsubpsace}, we deduce
$\Delta s_L = U_L^\top \Delta s_\gamma = O(\Xi^{-1} \epstol)$ and
$\Delta s_S = U_S^\top \Delta s_\gamma = O(\epstol)$.
Using \eqref{eq:cond:boundinvariantsubspace} and \eqref{eq:cond:inversecondensed},
the error in the large space $\Delta s_L$ annihilates in the backsolve:
\begin{equation}
  \label{eq:cond:boundhykkt1}
  K_\gamma^{-1} \Delta s_\gamma = U_L \Sigma_L^{-1} \Delta s_L + U_S \Sigma_S^{-1} \Delta s_S  = O(\epstol)
  \; .
\end{equation}
Finally, using \eqref{eq:cond:invperturbed}, we get
\begin{equation}
  \widehat{G} \widehat{K}_\gamma^{-1} \Delta s_\gamma \approx
  \widehat{G} (I - K_\gamma^{-1}\Delta K_\gamma) K_\gamma^{-1} \Delta s_\gamma \; .
\end{equation}
Using \eqref{eq:cond:boundhykkt1}, the first term is $\widehat{G} K_\gamma^{-1} \Delta s_\gamma = O(\epstol)$.
We have in addition
\begin{equation}
  G K_\gamma^{-1}\Delta K_\gamma (K_\gamma^{-1} \Delta s_\gamma)  =
  \big[ G U_L \Sigma_L^{-1} \Gamma_L + G U_S \Sigma_S^{-1} \Gamma_S \big] (K_\gamma^{-1} \Delta s_\gamma) \; .
\end{equation}
Using again \eqref{eq:cond:invariantsubpsace}:
$G U_L = G Y + O(\Xi)$ and $G U_S = O(\Xi)$. Hence
$G U_L \Sigma_L^{-1} \Gamma_L = O(1)$ and $G U_S \Sigma_S^{-1} \Gamma_S = O(1)$.
Using \eqref{eq:cond:rhserrorfull}, we have $K_\gamma^{-1} \Delta G^\top = O(\epstol)$,
implying $\Delta G K_\gamma^{-1} \Delta K_\gamma (K_\gamma^{-1} \Delta s_\gamma) = O(\Xi^{-1} \epstol^{2})$.
Assumption~\ref{hyp:cond:wellcond} implies that $\Xi^{-1} \epstol^2 \ll \epstol$,
proving \eqref{eq:cond:errorrgamma}.
\end{proof}

\end{document}
